\section{Implementation inventory and evidence}
\label{app:inventory}
This appendix records source deductions at implementation commit
\texttt{56a02ab}; analysis files are additional worktree changes. It is a static
protocol audit, not a build, concurrency test, or assertion that every path
is correct. The inventory is determined by
\src{src/command_parser/lock_target.rs}:110--154,183--215 and
\src{crates/libdlock/src/dlock2.rs}:55--81, rather than older README lists.
There are 21 DLock2 target variants, including the two PQ backends and both
Rust/C shuffle and CFL implementations. Similar names do not establish
identical execution, timing, or fairness contracts.
This is a shared factual inventory, not a combined analytical model.
Performance conclusions are in Section~\ref{sec:performance}; fairness
conclusions are in Section~\ref{sec:fairness}, each with its own comparison.

For compact source references below, \src{d2/} means
\src{crates/libdlock/src/dlock2/}, \src{d1/} means
\src{crates/libdlock/src/dlock/}, and \src{lib/} means
\src{crates/libdlock/src/}. Paths beginning \src{c/} or \src{src/} are
relative to the repository root. Line ranges refer to this baseline.

\subsection{All DLock2 targets}
\small
\begin{longtable}{@{}P{0.14\textwidth}P{0.51\textwidth}P{0.27\textwidth}@{}}
\toprule
Target enum & Model instance, budget, and guarantee boundary & Source evidence \\
\midrule
\endfirsthead
\toprule Target enum & Model instance, budget, and guarantee boundary & Source evidence \\
\midrule
\endhead
\bottomrule
\endfoot
FC & Inactive nodes publish through a head CAS; each pass scans the retained list to its end. No count/time cap. Cleanups run at pass multiples of 500 and remove sufficiently old nodes. Demand is list visits per completion, with election/activation retries separate; scan order does not establish usage fairness. & \src{d2/fc/lock.rs}:17,53--138,152--190. \\
\addlinespace
FCBan & Full FC-style scan plus combiner-side deadline tests. A skip does not complete a request. Penalties use elapsed work intervals and a counter of retained registrations, not instantaneous pending clients; eligibility and exposed server idle must be separated. & \src{d2/fc_ban/lock.rs}:59--62,103--138,148--167,182--205. \\
\addlinespace
CC & One tail swap per call, input/link publication in the predecessor slot, completion/wait flags, and baton handoff. At most 64 executed delegates per batch, possibly fewer on missing links. FIFO request order is not equal usage. & \src{d2/cc/lock.rs}:23,52--86,97--121. \\
\addlinespace
CCBan & TSC/backoff admission before enqueue, then queue combining. Cap 16 rather than CC's 64. Charge uses work intervals times a registration counter, which increments on first thread-local initialization. Thus the CC--CCBan delta is not a pure fairness-policy ablation. & \src{d2/cc_ban/lock.rs}:50--96,145--184. \\
\addlinespace
DSM & Per-client two-node toggle, tail enqueue, linked FIFO service. The post-execution test \texttt{counter > H} with default 64 permits 65 callbacks, subject to earlier successor/lookahead exits. It is not global double-buffer rotation. Unpublished successors can delay handoff. & \src{d2/dsm/lock.rs}:27,69--101,112--176. \\
\addlinespace
FcSL & Callers insert into a shared ordered SkipSet; a pass tries at most 64 pops and checks completion before executing. Successful service deactivates the node; a new call republishes. Charges include time since the previous timestamp, not only the delegate. Newcomer mean-charge initialization differs from a cumulative baseline. & \src{d2/fc_sl/lock.rs}:72--99,102--150,161--195. \\
\addlinespace
FcPqBTree & Shared sequential BTreeSet under combiner ownership, with a 64-slot announcement ring and at most 64 PQ pops. Ordinary queue work scales with arrivals/pops and queue size; allocation and migrating metadata matter. Exact-min proof is conditional on fixed continuous selectability and the billed metric. & \src{d2/fc_pq/lock.rs}:109--255,275--313; \src{lib/sequential_priority_queue.rs}:39--61. \\
\addlinespace
FcPqBHeap & Same admission, service policy and buffers as the tree version; BinaryHeap reverses ordering to pop minimum usage. Heap growth can relocate storage. Backend locality costs cannot be inferred from the identical min-credit policy. & \src{d2/fc_pq/lock.rs}:145--255; \src{lib/sequential_priority_queue.rs}:14--37. \\
\addlinespace
SpinLock & Contended acquisition operates on one AtomicBool with backoff; release stores false. The wrapper executes one delegate on the caller. Count attempts separately from ownership transfers; no FIFO or usage guarantee is supplied. & \src{lib/spin_lock.rs}:30--50; \src{d2/spinlock.rs}:41--46. \\
\addlinespace
MCS & Tail swap, predecessor link publication, own-node polling, and successor flag handoff or empty-tail CAS. One caller-executed operation per epoch. Constant successful queue operations do not bound a delayed link publication or protected-data migration latency. & \src{d2/mcs.rs}:78--107,124--152; \src{d2/spinlock.rs}:41--46. \\
\addlinespace
CLH & Tail swap, predecessor-flag polling, own-node release, and predecessor-node recycling. FIFO acquisition, not equal service. Polling is distributed over nodes but enqueue is a shared hotspot. & \src{d2/clh.rs}:98--124,171--185. \\
\addlinespace
Ticket & Fetch-add ticket admission and shared \texttt{now\_serving} polling; release advances the turn. Logical sharer invalidation demand is distinct from message count and latency. FIFO turn order does not bound heterogeneous service spread. & \src{d2/ticket.rs}:37--53. \\
\addlinespace
Mutex & One std::sync::Mutex guard and one delegate. The adapter supplies no fairness policy and does not pin a standard-library/kernel waiting implementation. Its opaque costs require the actual runtime and OS, not a guessed futex constant. & \src{d2/mutex.rs}:21--24,36--44. \\
\addlinespace
PthreadMutex & Explicit PTHREAD\_MUTEX\_NORMAL, lock/delegate/unlock. Libc/kernel contention, wakeup and scheduling remain external parameters; the type name does not supply an admission fairness bound. & \src{d2/pthread_mutex.rs}:60--65,89--100. \\
\addlinespace
ShflLock & Caller execution with fast-path ownership, queue promotion and same-NUMA shuffling. Leader traversal can stop at queue boundaries or a PRNG condition; no fixed request batch. \texttt{wcount} is not a delegate count budget. & \src{d2/shfl_lock/lock.rs}:23--30,200--289,298--417. \\
\addlinespace
ShflLockC & RawCAqs FFI reaches the Cargo-built AQS source. Same placement-policy family, but count its own traversal, promotion and lock-byte events; one delegate per acquisition, not combining. & \src{d2/c_aqs.rs}:75--88; \src{c/shfllock/aqs.c}:122--184,220--313. \\
\addlinespace
CFL & Queue shuffling plus global per-core/per-NUMA runtime accounting; a node-runtime scan and variable queue traversal affect policy demand. The fairness entity is not automatically an FC-PQ requester. Executor topology changes both policy work and $M_H$. & \src{d2/cfl.rs}:354--390,397--534,554--684,721--745. \\
\addlinespace
CFLC & RawCCfl FFI reaches the C implementation's runtime-based queue reshaping. One delegate per hold. Global core/node counters and a 100,000-tick node-spread threshold are source policy inputs, not a per-thread $C_{\max}$ theorem. Rust CFL also uses core/node accounting; neither should be described as an ideal per-request virtual-time scheduler. & \src{d2/c_cfl.rs}:75--90; \src{c/cfl/cfl.c}:187--304,312--445. \\
\addlinespace
FcC & Full active-list scan; completion is the delegate-pointer release to null after the response write. Cleanup uses 50-pass aging rather than Rust FC's 500. Waiters yield/retry after 100 pause iterations. The active scan has no enforced time cap merely because the build defines a timing macro. & \src{lib/c_binding/flatcombining.rs}:75--131; \src{c/FlatCombining/original/flatcombining.c}:13--56,90--129. \\
\addlinespace
CcC & FIFO tail queue, predecessor wait/completed flags, at most 64 callbacks while a successor exists, and a trailing-node combiner baton. This cap is established in the C source, not transferred from the Rust name. & \src{lib/c_binding/ccsynch.rs}:78--134; \src{c/CCsynch/ccsynch.c}:39--88. \\
\addlinespace
USCL & One caller-held critical section with slice admission, spin/yield, futex wait/wake and banned-time sleep. The nominal 2\,ms slice uses compile-time cycle conversion. Charge is held TSC ticks times an integer weight quotient, not delegate-only time. Default lazy C weighting is not an explicit DLock2 requester-weight policy. & \src{d2/uscl.rs}:21--41; \src{c/u-scl/fairlock.h}:91--185,188--305. \\
\end{longtable}
\normalsize

\paragraph{The actual C build.}
\src{crates/libdlock/build.rs}:18--33 compiles
\src{c/CCsynch/ccsynch.c},
\src{c/FlatCombining/original/flatcombining.c},
\src{c/u-scl/fairlock.c}, \src{c/shfllock/aqs.c}, and \src{c/cfl/cfl.c}.
\src{c/u-scl/fairlock.c}:1 includes the header containing its implementation.
The build defines \texttt{CYCLE\_PER\_US=2400}; this is not a measurement of
invariant-TSC frequency. A defined \texttt{FC\_THREAD\_MAX\_CYCLE} does not
establish a time bound in the original FC scan, which traverses to null.
Binding headers are in \src{crates/libdlock/binding/wrapper.h}:1--5.
The raw-lock DLock2 wrapper executes exactly one delegate between acquire and
release (\src{d2/spinlock.rs}:41--46), so these wrappers do not gain delegation
batching from sharing the DLock2 interface.

\subsection{FC-PQ proof and measurement obligations}
\begin{enumerate}
\item \textbf{Admission and visibility.}
\src{d2/fc_pq/lock/buffer.rs}:52--84 reserves by tail fetch-add, waits for
capacity, then publishes a valid flag. The iterator snapshots at most 64
tickets and spins on each valid flag before consuming it (:96--150).
This supports Proposition~\ref{prop:visibility}'s FIFO service-count bound,
not a wall-clock deadline. Counter wraparound and memory correctness are
excluded from that abstraction.
\item \textbf{Selectable versus merely pending.}
\src{d2/fc_pq/lock.rs}:175--255 buffers completed nodes, removes inactive ones,
and may reinsert reactivated nodes. A request becoming ready while absent
from the PQ need not be selected at the same boundary as a globally visible
minimum. Newcomer credit is the mean operation charge (:157--162).
No uniform admission-independent initial-spread bound follows.
\item \textbf{Attempted starvation prevention.}
The supported heap/tree backends pop a minimum
(\src{lib/sequential_priority_queue.rs}:14--61). Before the next queue mutation,
the post-pop clamp at \src{d2/fc_pq/lock.rs}:193--197 cannot decrease that
minimum. Its threshold cannot promote an entry that has not been popped.
\item \textbf{Billing and return.}
The outer timestamps bracket input access, delegate call and result write
(:209--217). The benchmark's inner useful-work timer has different boundaries
(\src{src/benchmark/dlock2/proportional_counter.rs}:81--110).
A combiner unlocks before checking its own completion (:286--295), so an
operation's finish and the caller's return are distinct.
\item \textbf{Misleading timer boundary.}
The same \texttt{begin} variable is set at pass entry and reset at every
delegate before \texttt{combiner\_time\_stat += end-begin}
(:129--131,209,258--263). This does not accumulate the entire pass when it
serves requests. FC-SL similarly resets its timestamp after each service
(\src{d2/fc_sl/lock.rs}:102--150). These statistics cannot be substituted for
full-pass occupancy without correcting the instrumentation.
\end{enumerate}
These are source-derived obligations and counterexamples, not silently fixed
production algorithms. In particular, proving an ideal policy's invariant
does not certify an unsafe concurrent implementation.

\subsection{FC-SL: concrete synchronization prerequisite}
\label{app:fcsl-race}
The DLock2 FC-SL newcomer path has a source-level data race. Its
\texttt{total\_usage} and \texttt{total\_served} fields are non-atomic
\texttt{SyncUnsafeCell<u64>} (\src{d2/fc_sl/lock.rs}:48,50).
Every \texttt{push\_node} reads the served counter, and zero-usage initialization
also reads total usage (:78--82). This path runs from \texttt{lock()} before
attempting the combiner mutex (:169--174), while another thread holding that
mutex can update both counters in \texttt{combine()} (:134--136).
There is no common synchronization protecting these conflicting accesses.
The per-request active/completion flags and a later SkipSet insertion do not
protect the global counters. Interior mutability and aligned x86 loads do
not legalize a Rust data race.

This is a correctness prerequisite for the concrete FC-SL instantiation,
not a claim about the magnitude of accounting error $E_c$. Moving initialization
under combiner ownership or adding synchronized counters with an explicit
snapshot-consistency contract are possible repairs, but no repair is made in
this analysis worktree. No executed race detector result or complete
concurrency audit is claimed; the conflicting source paths are sufficient
to identify this missing synchronization. Other variants have not thereby
been certified race-free.

\subsection{Legacy DLock1 coverage}
The 11 target variants are listed in
\src{src/command_parser/lock_target.rs}:41--66 and constructed at :83--105.
DLock1 passes callbacks over a guard; it is not the same API/charge boundary as
DLock2's function-delegate input/result path. The table accounts for every
legacy target, rather than transferring a DLock2 bound by matching its name.

\small
\begin{longtable}{@{}P{0.14\textwidth}P{0.51\textwidth}P{0.27\textwidth}@{}}
\toprule
Target & Instance and distinguishing source behavior & Evidence \\
\midrule
\endfirsthead
\toprule Target & Instance and distinguishing source behavior & Evidence \\
\midrule
\endhead
\bottomrule
\endfoot
FC & Callback-based full linked-list scan, with Parker wake/completion events; no fixed count/time cutoff in the scan. & \src{d1/fc/fclock.rs}:71--100. \\
\addlinespace
FCBan & Combiner-side banned-until filter over the retained list. Full scan and eligibility tests remain; do not transplant DLock2's exact penalty formula without checking its charge/update code. & \src{d1/fc_fair_ban.rs}:77--123,165--203. \\
\addlinespace
FCBanSlice & Despite its name and declared slice constant, the early-return slice test is commented out. The active combiner traverses to list end. Accumulated work affects an adaptive caller wait timeout, not an enforced pass deadline. & \src{d1/fc_fair_ban_slice.rs}:82--154,210--224. \\
\addlinespace
CC & Current combiner request is included in the \texttt{H=16} loop. Up to 16 callbacks before baton handoff; different from DLock2 CC's 64. & \src{d1/ccsynch.rs}:104--132. \\
\addlinespace
CCBan & Separately executes the current request, then up to 16 followers: up to 17 callbacks per episode. Adds pre-enqueue banned-time delay. Thus its total budget differs from DLock1 CC despite the same loop constant 16. & \src{d1/ccsynch_fair_ban.rs}:61--75,99--118,153--197. \\
\addlinespace
FCSLNaive & Usage-keyed concurrent SkipMap, not a different non-concurrent container inferred from ``Naive.'' Nominal accumulated work threshold 240,000 TSC units, checked after an item; callers retry after timed waiting without an explicit combiner baton. & \src{d1/fc_sl_naive.rs}:22--40,57--110,119--153. \\
\addlinespace
FCSL & Same SkipMap/threshold but adds completion tracking and wakes a remaining queued node as next combiner. Charging includes dequeue and surrounding work. Neither its post-item threshold nor calibrated cycle convention establishes preemptive execution. & \src{d1/fc_sl.rs}:22--40,57--126,134--166. \\
\addlinespace
RCL & Dedicated server thread scans and executes posted callbacks. Generic conversion intentionally returns None, but a separate benchmark path constructs it. Server-resource accounting is essential; this target is not absent from DLock1. & \src{src/benchmark/dlock.rs}:107--145; \src{d1/rcl/rclserver.rs}:85--98; \src{d1/rcl/rclthread.rs}:42--97. \\
\addlinespace
Mutex & Non-delegation OtherLocks baseline; uses the mutex callback adapter, not a parameterized Parker. Runtime/OS waiting remains a separate model component. & \src{src/command_parser/lock_target.rs}:98--102; \src{lib/dlock.rs}:55--119. \\
\addlinespace
SpinLock & Non-delegation OtherLocks baseline, not converted into a futex-Parker implementation by selecting the block waiter option. & \src{src/command_parser/lock_target.rs}:98--102; \src{src/benchmark/dlock.rs}:38--57. \\
\addlinespace
USCL & OtherLocks wrapper over the scheduler-cooperative acquisition protocol; no DLock1 Parker parameter is applied. Slice/OS costs must be modeled from its underlying implementation. & \src{src/command_parser/lock_target.rs}:98--102; \src{lib/u_scl.rs}:41--51. \\
\end{longtable}
\normalsize

\paragraph{Parking is a model dimension.}
For delegation targets, the default \texttt{All} waiter choice expands into
SpinParker and BlockParker variants
(\src{src/benchmark/dlock.rs}:38--57). SpinParker uses polling/backoff and may
yield (\src{lib/parker/spin_parker.rs}:26--63); BlockParker adds futex wait and
wake events (\src{lib/parker/block_parker.rs}:18--100). Equal service-credit
bounds do not remove scheduler wakeup delays. Baselines in OtherLocks do not
inherit this Parker selection.

\paragraph{RCL changes the resource budget, not only migration.}
\src{src/benchmark/dlock.rs}:125--145 starts the server using
\texttt{num\_cpu-1} and supplies that reduced CPU count to client work while
retaining the thread count. The server thread requests affinity at
\src{d1/rcl/rclthread.rs}:42--97. Assuming successful pinning, a single-server
run can eliminate rotating-executor migration after warmup, but still pays
request publication/completion and scanning. Its server is a resource in
Proposition~\ref{prop:resource}, and worker capacity is reduced; an apples-to-
apples comparison must hold total hardware resources fixed. This is source
reachability and requested placement, not a measurement that affinity always
succeeds. The exported DLock2 \src{d2/rcl.rs} modules have no DLock2 enum/CLI
variant (\src{lib/dlock2.rs}:25--26,55--81); direct use is not a dispatched
DLock2 benchmark and is not asserted working here.

\section{Reproduction and claim/evidence boundary}
\label{app:artifact}
The standalone manuscript is \src{analysis/logp/paper.tex}.
\src{performance.tex} and \src{fairness.tex} contain separate analyses;
\src{section.tex} supplies their common scope and final discussion.
The dependency-free artifact is \src{analysis/logp/check_model.py},
requiring Python 3.9 or newer. From the
worktree root:
\begin{quote}\ttfamily\small
python3 analysis/logp/check\_model.py\par
\quad --output .worktree/logp/verification.json
\end{quote}
The line break above is typographical: run it as one command. The adjacent
README supplies copyable commands for the checker and \texttt{latexmk} build.
Outputs are in the ignored \src{.worktree/logp/} directory. All numerical
illustrations in the section are exact normalized model values, not machine
measurements.

\begin{center}\small
\begin{tabular}{lrl}
\toprule
Check & Cases & Enumerated evidence \\
\midrule
\multicolumn{3}{l}{\emph{Performance}}\\
Crossover &5,120& Rational inequality cases \\
Batch frontier &5,832& Brute-force versus algebraic integer intervals \\
\addlinespace
\multicolumn{3}{l}{\emph{Fairness}}\\
Spread &44,400& One-step normalized integer invariant transitions \\
Accounting error &164,421& True/virtual-credit error transitions \\
Waiting bound &22,587& Competing transitions and terminal target choices \\
Charge bias &256& Complete equal-billed-service cycles \\
Announcement visibility &1,485& FIFO drain/service schedules \\
\addlinespace
\multicolumn{3}{l}{\emph{Cross-checks}}\\
Counterexamples &7& Source-inspired and metric witness families \\
\bottomrule
\end{tabular}
\end{center}
All eight checks passed in both standard and optimized (\texttt{python3 -O})
executions; the generated JSON reports were byte-identical. Their case counts
have different meanings and are not numbers of independent concurrent Rust
tests. Parameter domains are explicit in the script. It enumerates all
permitted choices within those domains rather than sampling random traces;
written proofs, not the finite domains, establish the conditional unbounded
results. The script uses explicit failures rather than optimization-removable
\texttt{assert} statements. Its JSON includes assumptions and example witnesses.

\begin{center}\small
\begin{tabular}{@{}P{0.22\textwidth}P{0.33\textwidth}P{0.36\textwidth}@{}}
\toprule
Claim & Evidence supplied & What is not inferred \\
\midrule
Abstract resource/migration bounds & Stated machine/residency assumptions and proofs & Guaranteed physical costs from empirical averages \\
Batch crossover and return frontier & Algebra, integer certificates, exact-rational checks & A hardware speedup or unconditional FC-PQ return deadline \\
Fixed-cohort service bounds & Inductive proofs, error reduction, finite transitions & Fairness under arbitrary admissions, suspension or credit resets \\
Charging bias & Exact fixed-cost derivation and complete-cycle checks & The magnitude of bias in actual benchmark data \\
FC-PQ visibility/newcomer/clamp limits & Source paths and abstract counterexamples & Executed Rust adversarial traces or full concurrency correctness \\
Coverage of targets & Enum/dispatch, budget and charge source audit & Successful compilation, safety, or performance of every variant \\
\bottomrule
\end{tabular}
\end{center}

The prescribed development environment previously failed evaluation because
\texttt{dotenv.resolved} had no defined value. The standard-library checks
and manuscript tooling use existing host programs instead; environment files
and lock implementations are unchanged. No performance suite is reported.
A full systems paper still needs independent calibration, corrected
measurement boundaries, controlled experimental resource accounting and
held-out prediction. These empirical requirements do not turn the proven
conditional statements into mere proposals, and the conditional statements
do not substitute for the empirical requirements.
